%! @@TITLE@@ — Unit @@UNITINT@@, Lesson @@LESSONID@@ — Group Activity
% GRADUAL RELEASE: this is the "you do together" block — 25 minutes, groups of 3–4.
% UNTIERED: every group works the same task (no Tier R/A/E). The vocabulary was built in the
% Guided Notes, so use it freely here. Reference implementations: unit01/lesson0{0,1}/activity.
\documentclass[10pt]{article}
\usepackage{@@PREFIX@@-article}
\usepackage{@@PREFIX@@-boxes}
% Coordinate grid: {xmin}{xmax}{ymin}{ymax}
\newcommand{\lgrid}[4]{%
  \draw[step=1, linegray!40, very thin] (#1,#3) grid (#2,#4);
  \draw[->, semithick, charcoal] (#1,0)--(#2,0) node[below right, font=\scriptsize]{$x$};
  \draw[->, semithick, charcoal] (0,#3)--(0,#4) node[left, font=\scriptsize]{$y$};}

\begin{document}

\pageheader{Unit @@UNITINT@@, Lesson @@LESSONID@@}{Group Activity}
% NAMESTRIP: no \namedateperiod here. The student writes their name once, on the cover.

\vspace{0.08in}

\begin{headlinebox}{forestbg}
\small\textbf{TODO Activity Title.} TODO: one motivating context in two sentences, then the charge
--- name the vocabulary from the notes the groups are expected to use, and finish with the standing
rule: \emph{for every answer, be ready to show me where you see it.}
\end{headlinebox}

\vspace{0.06in}

% FOUR parts, ~15–18 lettered sub-questions, 2–3 pages at 10pt. Parts 1–3 apply the notes; part 4
% is always MODEL IT — the lesson's modelling standard, worked in a real context.
\begin{scenariobox}[1. TODO: the straightforward case]{forest}
TODO: the scenario. Pre-draw every graph --- never ask students to sketch one from scratch.
\begin{enumerate}[label=\textbf{\alph*.}, leftmargin=1.8em, itemsep=5pt]
  \item TODO \blank{2.0cm}
  \item TODO --- an explain item.
        \writelines{2}
\end{enumerate}
\end{scenariobox}

\boxguard[22]
\begin{scenariobox}[2. TODO: the contrast case --- carries the crux question]{navy}
TODO: the same moves on a case that varies one condition. One item here is the \textbf{crux} ---
the one that surfaces the lesson's target misconception.
\begin{enumerate}[label=\textbf{\alph*.}, leftmargin=1.8em, itemsep=5pt]
  \item TODO \blank{2.0cm}
\end{enumerate}
\end{scenariobox}

\boxguard[16]
\begin{scenariobox}[3. TODO: put the two side by side]{forest}
TODO: the comparison the debrief will build on --- what the two cases share, what differs, and one
"write a sentence explaining why" item.
\begin{enumerate}[label=\textbf{\alph*.}, leftmargin=1.8em, itemsep=5pt]
  \item TODO \writelines{2}
\end{enumerate}
\end{scenariobox}

\boxguard[22]
\begin{scenariobox}[4. Model It --- TODO Title]{navy}
TODO: a fresh real context carrying the lesson's modelling standard. Write the model, solve it in a
\texttt{work} block, interpret the answer, then one "what if we change a number?" item that tests
the concept rather than the procedure.
\begin{enumerate}[label=\textbf{\alph*.}, leftmargin=1.8em, itemsep=5pt]
  \item TODO: what do the parameters mean here? \writelines{2}
  \item TODO: solve it.
    \begin{work}
      % TODO — identical in the key.
      40 - 8w &= 0 \\
            w &= 5
    \end{work}
    TODO: interpret. \writelines{2}
  \item TODO: the "what if?" concept check. \writelines{2}
\end{enumerate}
\end{scenariobox}

\end{document}
